\begin{document}
\maketitle

\begin{abstract}
On every finite set of disjoint sixteen-mode clusters we prove exactly
that the DET16 Hamiltonian
\[
 H_0=\sum_x(\id-P_{\phi,x})
\]
has the unique product ground state
$\bigotimes_x(|0\rangle_x+e^{i\phi}|F\rangle_x)/\sqrt2$, spectrum
$\{0,1,\ldots,|\Lambda|\}$, and gap one independently of volume.
The proof includes the graded-CAR argument: the determinant monomial has
even degree sixteen, so local projectors on disjoint clusters commute
without an ordering sign.  We also prove local Spin(10) Gauss invariance,
fermion parity, translation symmetry, absence of spontaneous symmetry
breaking, and absence of torus degeneracy.

For a bounded finite-range even perturbation, the
Michalakis--Zwolak stability theorem applies because frustration
freeness, Local-TQO, and the local gap are exact, with
$\Delta_0(\ell)=0$ and $\gamma(r)=1$.  It gives a
volume-independent threshold $t_*>0$ and perturbed gap at least $1/2$.
For finite $\mathbb Z_2/\mathbb Z_4$ dynamical links on an open chain,
the Fröhlich--Pizzo Lie--Schwinger theorem applies after an exact
ambient-cell and endpoint repair.  It gives a unique physical ground
state and gap at least $\Delta/2$, uniformly in chain length, for
$|t|<\Delta\tau_{\rm FP}/m$, where
$\Delta=\min(1,h_E)$ and $\tau_{\rm FP}>0$ is only partially explicit
in the cited proof.  It does \emph{not} give the proposed numerical
threshold $1/(2\cdot16z)$ for ordinary hopping.  The numeric twins pass
all checks and find no finite-volume closing in their declared scans.
The final box now isolates compact-group rotor links, wall mixing, and
the three-spatial-dimensional placement; none is included in the
finite-group chain theorem.
\end{abstract}

\tableofcontents

\section{Concrete model and frozen finite input}

\subsection{One DET16 cluster}

Let $K\cong\C^{16}$ carry the positive-chirality spin representation
$\rho_{16}$ of $\operatorname{Spin}(10)$ and let
$\F(K)=\bigoplus_{n=0}^{16}\bigwedge^nK$ be its fermionic Fock space.
For an ordered orthonormal basis $(e_1,\ldots,e_{16})$, write
\[
 \Omega^\dagger
 =a^\dagger(e_1)\cdots a^\dagger(e_{16}),\qquad
 |F\rangle=\Omega^\dagger|0\rangle ,
\]
and $\Omega=(\Omega^\dagger)^*$.  Put
\[
\begin{aligned}
 A_\phi
 &=P_0+P_{16}+e^{-i\phi}\Omega+e^{i\phi}\Omega^\dagger,\\
 |\Omega_\phi\rangle
 &=\frac{|0\rangle+e^{i\phi}|F\rangle}{\sqrt2},\\
 P_\phi&=\frac12A_\phi
 =|\Omega_\phi\rangle\langle\Omega_\phi|,\qquad
 q_\phi=\id-P_\phi.
\end{aligned}
\tag{1.1}\label{eq:cluster}
\]
The exact finite computation in
\texttt{verification/v1002\_spin10\_mirror\_projector.py} verifies
\[
 \Omega^\dagger\Omega=P_{16},\quad
 \Omega\Omega^\dagger=P_0,\quad
 P_\phi^2=P_\phi=P_\phi^*,\quad
 \rank P_\phi=1,
\tag{1.2}\label{eq:frozen}
\]
and all forty-five $\mathfrak{so}(10)$ commutators vanish.  These are
inputs to this memorandum; the algebraic consequences needed below
are also proved directly.

\subsection{Many disjoint clusters}

Let $\Lambda$ be any nonempty finite set of clusters.  Fix an ordering
only to identify the global Fock space with an ordinary tensor product,
\[
 \F\!\left(\bigoplus_{x\in\Lambda}K_x\right)
 \cong\widehat{\bigotimes}_{x\in\Lambda}\F(K_x),
\tag{1.3}\label{eq:gradedtensor}
\]
where the hat recalls that the intrinsic product is graded.  Operators
have fermion parity $|A|\in\{0,1\}$.  If $A_x$ and $B_y$ are homogeneous
and supported on disjoint clusters, the CAR give
\[
 A_xB_y=(-1)^{|A_x||B_y|}B_yA_x.
\tag{1.4}\label{eq:gradedcommute}
\]
The terms $P_{0,x}$ and $P_{16,x}$ preserve particle number, while
$\Omega_x$ and $\Omega_x^\dagger$ contain sixteen odd generators.
Every summand of $P_{\phi,x}$ is therefore even.  Hence
\[
 [P_{\phi,x},P_{\phi,y}]=[q_{\phi,x},q_{\phi,y}]=0
 \qquad(x\ne y)
\tag{1.5}\label{eq:commute}
\]
with no ordering-dependent sign.  Define
\[
 H_{0,\Lambda}=\sum_{x\in\Lambda}q_{\phi,x}.
\tag{1.6}\label{eq:H0}
\]
Unless stated otherwise, the same phase $\phi$ is used at every site.

\section{Exact spectrum and the volume-independent unit gap}

\begin{lemma}[Order independence of even local products]
\label{lem:order}
For any subset $S\subseteq\Lambda$, the product
$\prod_{x\in S}P_{\phi,x}$ is independent of the order chosen for its
factors.  Under any ordered tensor-product identification
\eqref{eq:gradedtensor}, it is the ordinary tensor product of the local
rank-one projectors on $S$ and the identity on $S^c$.
\end{lemma}

\begin{proof}
Each $P_{\phi,x}$ is even.  Equation \eqref{eq:gradedcommute} therefore
has sign $(-1)^{0\cdot0}=1$ for every transposition of two factors.
Any two orderings differ by transpositions, proving the first claim.
An even operator crosses all preceding graded tensor factors without a
Jordan--Wigner parity string.  Consequently its representative is the
ordinary local matrix on its own factor and identities elsewhere.
Multiplying these representatives proves the second claim.
\end{proof}

\begin{theorem}[T1: exact DET16 product spectrum]
\label{thm:T1}
Let $N=|\Lambda|$ and $d=\dim\F(\C^{16})=2^{16}$.  The commuting
Hamiltonian \eqref{eq:H0} has
\[
 \operatorname{spec}(H_{0,\Lambda})=\{0,1,\ldots,N\}.
\tag{2.1}\label{eq:spectrum}
\]
The multiplicity of $k$ is
\[
 \binom Nk(d-1)^k.
\tag{2.2}\label{eq:multiplicity}
\]
Its unique normalized ground state is
\[
 |\Omega_{\Lambda,\phi}\rangle
 =\widehat{\bigotimes}_{x\in\Lambda}|\Omega_\phi\rangle_x,
\tag{2.3}\label{eq:productground}
\]
and
\[
 \gap(H_{0,\Lambda})=1
\tag{2.4}\label{eq:gapone}
\]
for every finite volume and every boundary condition.
\end{theorem}

\begin{proof}
On one cluster,
$\F(K_x)=\operatorname{Ran}P_{\phi,x}\oplus
\operatorname{Ran}q_{\phi,x}$ with dimensions $1$ and $d-1$.
By Lemma~\ref{lem:order}, the global Hilbert space is the orthogonal
direct sum
\[
 \bigoplus_{S\subseteq\Lambda}
 \left(
 \bigotimes_{x\in S}\operatorname{Ran}q_{\phi,x}
 \otimes
 \bigotimes_{y\notin S}\operatorname{Ran}P_{\phi,y}
 \right).
\tag{2.5}\label{eq:syndrome}
\]
On the summand labelled by $S$, exactly the projectors $q_{\phi,x}$
with $x\in S$ act as one and all others act as zero.  Thus
$H_{0,\Lambda}$ acts as the scalar $|S|$.  This proves
\eqref{eq:spectrum}.  For $|S|=k$ there are $\binom Nk$ choices of
$S$, and every chosen excited factor has dimension $d-1$, proving
\eqref{eq:multiplicity}.

At energy zero, $S=\varnothing$ is the only summand and every local
factor lies in the one-dimensional range of $P_{\phi,x}$.  Its
normalized vector is exactly \eqref{eq:productground}; evenness makes
the product independent of the temporary cluster ordering.  The next
spectral value is one, proving \eqref{eq:gapone}.  No edge or torus
identification enters \eqref{eq:H0}, so the conclusion is independent
of boundary condition.
\end{proof}

\section{Gauge, parity, translations, and degeneracy}

\begin{lemma}[The chiral determinant character is trivial]
\label{lem:det}
For every $g\in\operatorname{Spin}(10)$,
\[
 \det\rho_{16}(g)=1.
\tag{3.1}\label{eq:detone}
\]
\end{lemma}

\begin{proof}
The map $\chi=\det\circ\rho_{16}:
\operatorname{Spin}(10)\to U(1)$ is a continuous group character.
Its differential at the identity is
\[
 d\chi(X)=\Tr(d\rho_{16}(X)).
\]
The trace vanishes on commutators.  Since
$\mathfrak{spin}(10)$ is semisimple,
$[\mathfrak{spin}(10),\mathfrak{spin}(10)]
=\mathfrak{spin}(10)$, so this linear functional vanishes identically.
Thus $\chi$ is locally constant.  The group $\operatorname{Spin}(10)$
is connected, and $\chi(e)=1$, hence $\chi(g)=1$ for every $g$.
\end{proof}

\begin{theorem}[T2: exact symmetries and no hidden degeneracy]
\label{thm:T2}
The Hamiltonian $H_{0,\Lambda}$ has the following properties.
\begin{enumerate}[label=\textup{(\alph*)}]
\item It is invariant under independent local Spin(10) transformations.
Equivalently, for every matter Gauss generator
$G_x^A=d\Gamma_x(T^A)$,
\[
 [H_{0,\Lambda},G_x^A]=0
 \quad(x\in\Lambda,\ A=1,\ldots,45).
\tag{3.2}\label{eq:gauss}
\]
\item It commutes with total fermion parity.  If $\phi$ is constant,
it is invariant under every lattice translation preserving $\Lambda$
and its boundary condition.
\item The finite-volume ground state is a symmetry singlet.  The
infinite-volume product state is the unique frustration-free ground
state, so none of these symmetries is spontaneously broken.
\item On a torus the ground-state degeneracy is exactly one.  In
particular, the model has no topological ground-state degeneracy.
\end{enumerate}
\end{theorem}

\begin{proof}
Second quantization fixes the vacuum.  On the filled vector it acts by
the top exterior power:
\[
 \Gamma(\rho_{16}(g))|F\rangle
 =(\det\rho_{16}(g))|F\rangle=|F\rangle
\]
by Lemma~\ref{lem:det}.  Hence $|\Omega_\phi\rangle$, $P_\phi$, and
$q_\phi$ are invariant.  This is true independently at every cluster,
so differentiation gives \eqref{eq:gauss}.  The forty-five zero
commutators in the frozen finite computation are the matrix twin of
this argument.

Both components of $|\Omega_\phi\rangle$ have even parity because
their particle numbers are $0$ and $16$.  Thus every $P_{\phi,x}$ and
$q_{\phi,x}$ is even and $H_0$ commutes with total parity.  A
translation merely permutes identical terms when $\phi$ is constant,
which proves translation invariance.

The product ground vector is fixed by all the stated symmetries.  For
the infinite system, let $\omega$ be any frustration-free state.
Then $\omega(q_{\phi,x})=0$ for every $x$.  Positivity and
$q_{\phi,x}\geq0$ imply that the one-site reduced state is supported
on $\ker q_{\phi,x}=\C|\Omega_\phi\rangle$.  A state with a pure
one-site marginal factorizes across that site; iterating over every
finite region gives
\[
 \omega|_A
 =\bigotimes_{x\in A}
 |\Omega_\phi\rangle\langle\Omega_\phi|
\]
for every finite $A$.  Local restrictions determine a quasi-local
state, so $\omega$ is the unique product state.  There are therefore
no distinct infinite-volume ground states related by a symmetry and
no spontaneous symmetry breaking.

Finally, Theorem~\ref{thm:T1} applies unchanged on a torus and gives a
one-dimensional ground space.  More concretely, all local constraints
$q_{\phi,x}=0$ force the same rank-one factor independently at every
site; no noncontractible cycle supplies an additional label.  This
proves the torus claim directly from the commuting-projector
structure.
\end{proof}

\section{The classical stability theorem, stated in full}

\subsection{Local-TQO and local gap}

For a finite region $B\subseteq\Lambda$, define
\[
 H_B=\sum_{x\in B}q_{\phi,x},\qquad
 P_B=\prod_{x\in B}P_{\phi,x}.
\tag{4.1}\label{eq:localHB}
\]
As an operator on the full system, $P_B$ is the identity outside
$B$.  If $A\subset B$ and $O_A$ is supported on $A$, then
\[
 P_BO_AP_B
 =\ip{\Omega_{A,\phi}}{O_A\Omega_{A,\phi}}P_B.
\tag{4.2}\label{eq:exactLTQO}
\]
Thus local ground states are exactly indistinguishable in the
interior: the Local-TQO error function is identically zero.  Theorem
\ref{thm:T1} applied to $B$ also gives
\[
 \gap(H_B)=1
\tag{4.3}\label{eq:localgap}
\]
for every nonempty $B$.

\begin{definition}[Finite-range perturbation strength]
\label{def:pert}
An even finite-range perturbation is a sum
\[
 V_\Lambda=\sum_{X\subseteq\Lambda}V_X,\qquad
 V_X=V_X^*,\qquad |V_X|=0,
\tag{4.4}\label{eq:pert}
\]
for which $V_X=0$ when $\operatorname{diam}X>R$ and
\[
 \sup_{x\in\Lambda}\sum_{X\ni x}\norm{V_X}\leq J_V
\tag{4.5}\label{eq:interactionnorm}
\]
uniformly in $\Lambda$.  Gauge covariance means that the complete
local term, including any prescribed parallel transporter, is
Spin(10)-invariant.  Evenness makes disjoint interaction terms commute
in the graded sense required by the local stability argument.
\end{definition}

\begin{theorem}[Michalakis--Zwolak stability theorem, full form used
here]
\label{thm:MZ}
Let $H_0$ be a finite-range frustration-free Hamiltonian on finite
subsets of $\mathbb Z^D$, with a volume-independent spectral gap
$\gamma>0$.  For every ball $b_u(r)$ let $P_{b_u(r)}$ be the ground
projection of the restricted Hamiltonian.  Assume:
\begin{enumerate}[label=\textup{(\roman*)}]
\item \emph{Local-TQO.}  For every ball $A=b_u(r)$, every bounded
operator $O_A$ supported in $A$, and
$A(\ell)=b_u(r+\ell)$,
\[
 \left\|
 P_{A(\ell)}O_AP_{A(\ell)}
 -\frac{\Tr(P_{A(\ell)}O_A)}
 {\Tr P_{A(\ell)}}P_{A(\ell)}
 \right\|
 \leq\norm{O_A}\Delta_0(\ell),
\tag{4.6}\label{eq:MZLTQO}
\]
where $\Delta_0(\ell)$ decays sufficiently rapidly and the admissible
topological-order length grows as a positive power of the system
diameter.
\item \emph{Local gap.}  Every restricted Hamiltonian on $b_u(r)$ has
no spectrum in $(0,\gamma(r))$, where $\gamma(r)>0$ decays at most
polynomially in $r$.
\item \emph{Quasi-local perturbation.}  The perturbation is a
$(J,f_0)$ interaction: it is a sum of bounded terms centred at lattice
sites whose range-$r$ norms are bounded by $Jf_0(r)$, with $f_0$
decaying rapidly enough for a Lieb--Robinson bound.
\end{enumerate}
Then there exist constants $J_0>0$ and $L_0\geq2$, independent of the
finite volume, such that whenever $J\leq J_0$ and the system diameter
$L\geq L_0$,
\[
 \gap(H_0+V)\geq\frac{\gamma}{2}.
\tag{4.7}\label{eq:MZgap}
\]
In the proof, after quasi-local spectral-flow block diagonalization,
the transformed interaction $W$ obeys
\[
 |\ip{\psi}{W\psi}|
 \leq cJ\,\ip{\psi}{H_0\psi},
\tag{4.8}\label{eq:MZrelative}
\]
where one may take
\[
 c=C_D\sum_{k=1}^{M}
 \frac{r_k^D}{\gamma(r_k)}
 \sum_{r_{k-1}<r\leq r_k}w(r),
\qquad
 0=r_0<r_1\leq\cdots\leq r_M=L,
\tag{4.9}\label{eq:MZc}
\]
$C_D$ is the volume of the unit ball in $\mathbb Z^D$, and $w$ is the
rapid-decay function of the transformed interaction.  The final
bootstrap may choose
\[
 J_0=\frac1{3c}.
\tag{4.10}\label{eq:MZJ0}
\]
The constants $w,c,J_0$ depend on the interaction-decay and
spectral-flow estimates, but not on volume.
\end{theorem}

This is Theorem~1 and Proposition~2 of
S.~Michalakis and J.~P.~Zwolak,
\emph{Stability of Frustration-Free Hamiltonians},
Communications in Mathematical Physics \textbf{322} (2013), 277--302.
It extends the commuting-projector theorem of Bravyi, Hastings, and
Michalakis.  We have included the hypotheses, conclusion, and the
relative-bound constant used in their proof so that no numerical
threshold is hidden.

